\documentclass[10pt,a4paper]{amsart}
\usepackage[margin=19mm]{geometry}
\usepackage{amsmath,amssymb,mathtools}
\usepackage[scr=rsfs,cal=euler]{mathalfa}
\usepackage{xfrac}
\usepackage[unicode,hidelinks,bookmarks=false]{hyperref}
\newcommand{\ie}{i.e.\ }
\newcommand{\cf}{cf.\ }
\newcommand{\resp}{respectively\ }
\newcommand{\Subsection}{\subsection}
\providecommand{\nobreakdash}{\nobreak\hbox{-}}
\setlength{\emergencystretch}{2em}
\multlinegap0pt
\usepackage{bm}
\usepackage{bbm}
\usepackage{dsfont}
\usepackage{xspace}
\usepackage{cancel}
\usepackage{mathtools}
\usepackage{amsthm}
\makeatletter
\@ifundefined{claim}{\newtheorem{claim}{Claim}}{}
\makeatother
\makeatletter
\def\P{\mathbb P}
\expandafter\ifx\csname named\endcsname\relax
\newenvironment{named}[1]
{\renewcommand{\thistheoremname}{#1}%
	\begin{genericthm}}
\fi
\expandafter\ifx\csname proofclaim\endcsname\relax
\newenvironment{proofclaim}[1][]%
{\noindent \emph{Proof.} {}{#1}{}}{\hfill
	$\Diamond$\vspace{1em}}
\fi
\makeatother
\title[When are random regular triangle-free graphs bipartite?]{When are random regular triangle-free graphs bipartite?}
\author{Gregory DeCamillis, Pu Gao}
\date{Source: arXiv:2609.17958v1 (CC BY 4.0)}
\begin{document}
\maketitle

\noindent\emph{Source and licence.} When are random regular triangle-free graphs bipartite?. Version arXiv:2609.17958v1 (CC BY 4.0), distributed under the Creative Commons Attribution 4.0 International licence, as recorded in the arXiv licence metadata for this exact version and shown on the abstract page https://arxiv.org/abs/2609.17958v1. Full rights notice: licenses/arxiv-2609.17958-CC-BY-4.0.txt.\par\smallskip

\noindent\textit{Editorial scope.} Counted unit: the Claim whose source label is \texttt{clm: m to p} in the article (source0.tex lines 1736--1741, source label \texttt{clm: m to p}), delivered together with the article's own complete original proof of it (source0.tex lines 1742--1774, one \texttt{proof} environment of 33 lines). No mathematics is written, repaired or re-derived: statement and proof are the article's own text line for line. Required context retained uncounted: none. Every definition, display and label of those lines is retained unchanged. Omitted from this package and not redistributed: 1--1735, 1775--3082. Nothing inside a retained excerpt is omitted or altered: every retained body is delivered exactly as the article's lines are written, so no omission receipt of any kind accompanies this package. Citations: the retained excerpts carry no citation command at all, so no bibliography entry of the article is transcribed and no reference is redistributed. Reference adaptation: none was needed and none was made; every reference inside the delivered unit resolves inside it, so no unresolved reference or citation is produced. Printed slips kept literal: no printed word, symbol, hypothesis or number of the counted statement or of its proof was altered. Layout and class: the article's own class is \texttt{article} with packages including amsthm, amssymb, amsmath, enumerate, comment, thm-restate, centernot, url, subfigure, enumitem, mathrsfs, xcolor, color, ulem; the wrapper is the lane's pinned \texttt{amsart} editorial wrapper at 10pt on A4, which restates the theorem-like environments the retained text needs and adds the article's own macro definitions that the retained text needs (\texttt{\textbackslash P}), with extra wrapper packages (bm, bbm, dsfont, xspace, cancel, mathtools). The delivered numbering is the unit's own. Assets: no retained span contains a figure environment, a graphics-inclusion command, a PDF-page inclusion, a table or a TikZ picture; no image file is copied into this package, the package declares no asset dependency and no image file is redistributed with it. Licence: version arXiv:2609.17958v1 of this article is distributed under the Creative Commons Attribution 4.0 International licence, as recorded in the arXiv licence metadata for this exact version and shown on the abstract page https://arxiv.org/abs/2609.17958v1. A full rights notice is delivered at licenses/arxiv-2609.17958-CC-BY-4.0.txt. Characters: every retained excerpt is pure ASCII, so the delivered engine, XeLaTeX with its default Latin Modern text font, reports no missing character.
\begin{claim}\label{clm: m to p}
    For $C \in \{1,A,B\}$,
    \[
    \P_{H_C'}[X_C = 0] =\Omega(\P_{H_C}[X_C = 0]).
    \]
\end{claim}

\begin{proof}
    Let $(U,V)$ be the bipartition of $H_C$. Since the number of edges in $H'_C$ is a binomial random variable with expectation $m'_C$, we have 
    \[
    \P_{H'_C}[|E(H'_C)|\leq m'_C] = \Omega(1).
    \] 
    Thus, 
    \begin{align*}
        \P_{H'_C}[X_C = 0] &\geq \P_{H'_C}[X_C = 0 \mathrel ||E(H'_C)| \leq m'_C] \cdot\P_{H'_C}[|E(H'_C)| \leq m'_C]\\
        &= \Omega(\P_{H'_C}[X_C = 0 \mathrel ||E(H'_C)| \leq m'_C]).
    \end{align*}
We conclude this proof by showing that for any $r \geq 1$ where $\P[|E(H'_C)| = r] \neq 0$, 
    \[
        \P_{H'_C}[X_C = 0 \mathrel | |E(H'_C)| = r] \leq \P_{H'_C}[X_C = 0 \mathrel||E(H'_C)| = r-1],
    \]
    since it would follow that 
    \begin{align*}
        \P_{H'_C}[X_C = 0\mathrel | |E(H'_C)| \leq m'_C] &= \Omega(\P_{H'_C}[X_C = 0\mathrel | |E(H'_C)| = m'_C])\\
        &= \Omega(\P_{H_C}[X_C = 0]).
    \end{align*}
    For any such $r$ (or $r = 0$), let $\mathcal G_r$ denote the total number of bipartite graphs on bipartition $(U,V)$ with $r$ edges and let $\mathcal T_r \subseteq \mathcal G_r$ be the set of graphs $G \in \mathcal G_r$ where $G \cup T_A \cup T_B$ is triangle-free. Let $t_r = |\mathcal T_r|$ and $g_r = |\mathcal G_r|$. It suffices to show that $t_r/g_r \leq t_{r-1}/g_{r-1}$. That is,
    \[
    \frac{t_{r}}{t_{r-1}} \leq \frac{g_{r}}{g_{r-1}} = \frac{\binom{|U||V|}{r}}{\binom{|U||V|}{r-1}} = \frac{|U||V| - r + 1}{r}.
    \]
    We prove this using a simple double counting argument. Consider a relation $\sim $ on $\mathcal T_r \times \mathcal T_{r-1}$ where $G \sim G'$ if $G' = G -e$ for some $e \in E(G)$. Let $\mathscr G$ be the bipartite graph with bipartition $(T_{r} , \mathcal T_{r-1})$ and edges defined by $\sim$. Let $d_r$ denote the minimum degree of graphs in $\mathcal T_{r}$ and $d_{r-1}$ denote the maximum degree of graphs in $\mathcal T_{r-1}$ in $\mathscr G$. Then, $d_rt_r \leq d_{r-1}t_{r-1}$, so it suffices to show that 
    \[
    \frac{d_{r-1}}{d_{r}} \leq \frac{|U||V| - r+1}{r}.
    \]
    Note that deleting any edge of a graph in $\mathcal T_{r}$ will result in a graph in $\mathcal T'_{r-1}$, so we have 
    \[
    d_r \geq r.
    \]
    Also, for any graph $T \in \mathcal T'_{H, V, m}$, at most $|U||V| - r + 1$ edges can be added to $T$ to obtain a bipartite graph with bipartition $(U, V)$. That is,  $d_{r-1} \leq |U||V| - r+1$. This completes the proof.
\end{proof}

\end{document}
